% AA-53 — pierścienie niezmienników.
\chapter{Pierścienie niezmienników grup skończonych}
\label{aa:ch:finite-invariants}
W rozdziale~\ref{aa:ch:binary-polyhedral} rozróżniliśmy
niezmienniki i półniezmienniki. Teraz wyznaczymy
całe pierścienie niezmienników dla dwóch nieskończonych
rodzin grup oraz metodę obliczania wymiarów ich części
jednorodnych. Pracujemy nad ciałem $\mathbb C$.

\section{Działanie na wielomianach}

\begin{definicja}[Pierścień niezmienników i szereg Hilberta]
\label{aa:def:invariants-hilbert}
Niech $V$ będzie skończenie wymiarową zespoloną
przestrzenią liniową, a $G\le\operatorname{GL}(V)$
grupą skończoną. Niech $S=\mathbb C[V]$ oznacza
pierścień funkcji wielomianowych na $V$,
a $S_d$ jego część jednorodną stopnia $d\ge0$.
Dla $g\in G$, $f\in S$ i $v\in V$ definiujemy
$(g\cdot f)(v)=f(g^{-1}v)$. Jest to działanie grupy,
gdyż $1\cdot f=f$ oraz
$(g\cdot(h\cdot f))(v)=f(h^{-1}g^{-1}v)
=((gh)\cdot f)(v)$.
Pierścieniem niezmienników nazywamy
$S^G=\{f\in S:g\cdot f=f\text{ dla każdego }g\in G\}$.
Jego szeregiem Hilberta nazywamy formalny szereg
$H_G(z)=\sum_{d\ge0}\dim_{\mathbb C}(S_d^G)z^d$.
\end{definicja}

Działanie zachowuje stopień oraz iloczyn, więc
$S^G=\bigoplus_{d\ge0}S_d^G$ jest podpierścieniem
z gradacją. Szereg Hilberta opisuje wymiary, ale sam
nie określa mnożenia ani prezentacji pierścienia.

\begin{lemat}[Operator Reynoldsa]
\label{aa:lem:invariants-reynolds}
Dla $G,V,S$ z Definicji~\ref{aa:def:invariants-hilbert}
odwzorowanie $\mathcal R:S\to S$ określone wzorem
$\mathcal R(f)=|G|^{-1}\sum_{g\in G}g\cdot f$
jest rzutem liniowym na $S^G$, zachowującym stopnie.
Dla $a\in S^G$, $f\in S$ zachodzi
$\mathcal R(af)=a\mathcal R(f)$.
\end{lemat}
\begin{proof}
Dla $h\in G$ lewostronne mnożenie przez $h$ permutuje
składniki sumy, więc $h\cdot\mathcal R(f)=\mathcal R(f)$.
Jeśli $f\in S^G$, każdy składnik sumy jest równy $f$.
Stąd obraz jest równy $S^G$ i $\mathcal R^2=\mathcal R$.
Z liniowości działania na $V$ wynika zachowanie stopni.
Ponadto $g\cdot(af)=a(g\cdot f)$, co po uśrednieniu
daje ostatnią równość.
\end{proof}

\begin{twierdzenie}[Skończona generowalność niezmienników]
\label{aa:tw:invariants-finite}
Niech $V$ będzie skończenie wymiarową zespoloną
przestrzenią liniową, a $G\le\operatorname{GL}(V)$
grupą skończoną. Wówczas pierścień niezmienników
$\mathbb C[V]^G$ jest skończenie generowaną algebrą
nad $\mathbb C$.
\end{twierdzenie}
\begin{proof}
Oznaczmy $S=\mathbb C[V]$. Niech $I$ będzie ideałem
w $S$ generowanym przez
wszystkie jednorodne niezmienniki dodatniego stopnia.
Z twierdzenia Hilberta o bazie pierścień $S$ jest
noetherowski. Istnieje więc skończona rodzina
jednorodnych niezmienników $f_1,\ldots,f_r$ dodatniego
stopnia generująca $I$ jako ideał $S$:
każdy z dowolnych skończonych generatorów $I$
jest skończoną kombinacją pierwotnych generatorów,
wystarczy zebrać wszystkie występujące w niej niezmienniki.
Jeśli $I=0$, pierścień niezmienników składa się ze stałych.

Dla jednorodnego niezmiennika $f$ stopnia $d>0$
mamy $f=\sum_i a_i f_i$. Biorąc składniki jednorodne
stopnia $d$, możemy przyjąć $a_i$ jednorodne stopnia
$d-\deg f_i$ (lub zerowe przy ujemnym stopniu).
Operator Reynoldsa daje $f=\sum_i\mathcal R(a_i)f_i$.
Każdy niezerowy $\mathcal R(a_i)$ ma stopień mniejszy
od $d$. Indukcja względem $d$, rozpoczynająca się od
stałych, daje $f\in\mathbb C[f_1,\ldots,f_r]$.
Rozkład na części jednorodne kończy dowód.
\end{proof}

\section{Wzór Moliena}

\begin{twierdzenie}[Wzór Moliena]
\label{aa:tw:molien}
Niech $V$ będzie skończenie wymiarową zespoloną
przestrzenią liniową, a $G\le\operatorname{GL}(V)$
grupą skończoną. Niech $H_G(z)$ będzie szeregiem
Hilberta pierścienia $\mathbb C[V]^G$ dla gradacji
według stopnia wielomianów. W pierścieniu
$\mathbb C[[z]]$ zachodzi
\[
 H_G(z)=\frac{1}{|G|}\sum_{g\in G}
 \frac{1}{\det(I-zg)}.
\]
\end{twierdzenie}
\begin{proof}
Oznaczmy przez $S_d$ przestrzeń wielomianów
jednorodnych stopnia $d$ na $V$ i połóżmy
$m=\dim_{\mathbb C}V$. Ograniczenie operatora Reynoldsa
do $S_d$ jest rzutem, którego ślad wynosi $\dim S_d^G$:
w bazie dostosowanej do obrazu i jądra rzut ma na
przekątnej jedynki na obrazie i zera na jądrze.
Jeśli wartości własne $g^{-1}$ na $V$, z uwzględnieniem
krotności, to $\lambda_1,\ldots,\lambda_m$,
w bazie funkcjonałów liniowych sprzężonej do bazy
wektorów własnych $g^{-1}$ monom stopnia $d$,
o wykładnikach $\alpha_1,\ldots,\alpha_m$,
ma wartość własną $\lambda_1^{\alpha_1}
\cdots\lambda_m^{\alpha_m}$ dla działania $g$.
Taka baza istnieje, bo $g$ ma skończony rząd
i jego wielomian minimalny dzieli wielomian bez
pierwiastków wielokrotnych $T^{|G|}-1$.
Sumując ślady po stopniach, otrzymujemy
\[
 \sum_{d\ge0}\operatorname{tr}(g|_{S_d})z^d
 =\prod_{j=1}^m(1-\lambda_jz)^{-1}
 =\det(I-zg^{-1})^{-1}.
\]
Iloczyny są formalnymi szeregami geometrycznymi;
w każdym stopniu występuje skończenie wiele składników.
Uśrednienie po $G$ i zastąpienie $g^{-1}$ przez $g$
w skończonej sumie daje wzór.
\end{proof}

Dla $G\le\operatorname{SL}_2(\mathbb C)$ mianownik
ma postać $1-\operatorname{tr}(g)z+z^2$.
W szczególności wystarczy znać liczności elementów
o poszczególnych śladach, aby wyznaczyć szereg Hilberta.

\section{Grupy cykliczne}

\begin{twierdzenie}[Niezmienniki grupy cyklicznej]
\label{aa:tw:cyclic-invariant-ring}
Niech $m\ge1$ będzie całkowita, $\zeta$ pierwotnym
pierwiastkiem $m$-tego stopnia z jedności, a
$G=\langle\operatorname{diag}(\zeta,\zeta^{-1})\rangle$.
Dla działania na $\mathbb C[x,y]$ jak w
Definicji~\ref{aa:def:invariants-hilbert} mamy
\[
 \mathbb C[x,y]^G\simeq
 \mathbb C[U,V,W]/(UV-W^m),
 \quad (U,V,W)\longmapsto(x^m,y^m,xy).
\]
Stopnie generatorów to $m,m,2$ i
\[
 H_G(z)=\frac{1+z^m}{(1-z^m)(1-z^2)}.
\]
\end{twierdzenie}
\begin{proof}
Monom $x^a y^b$ jest niezmiennikiem dokładnie wtedy,
gdy $m$ dzieli $a-b$. Ponieważ monomy tworzą bazę
wektorów własnych generatora, każdy niezmiennik
jest kombinacją takich monomów. Jeśli $a\ge b$,
monom ma postać $(xy)^b(x^m)^{(a-b)/m}$;
dla $a<b$ postępujemy analogicznie z $y^m$.
To dowodzi surjektywności wskazanego homomorfizmu.

Równość $UV=W^m$ pozwala zredukować każdy monom
w ilorazie do jednego z monomów $W^k$,
$W^kU^j$, $W^kV^j$ ($k\ge0$, $j\ge1$).
Ich obrazy w $\mathbb C[x,y]$ są parami różnymi
monomami (rozróżnia je znak i wartość $a-b$
oraz $\min(a,b)$). Są zatem liniowo niezależne.
Dowodzi to iniektywności, a zliczanie ich stopni daje
\[
 H_G(z)=\frac{1}{1-z^2}
 \left(1+2\sum_{j\ge1}z^{mj}\right)
 =\frac{1+z^m}{(1-z^2)(1-z^m)}.
\]
\end{proof}

\begin{przyklad}[Grupa rzędu dwa]
Dla $m=2$ niezmiennikami są dokładnie wielomiany
o monomach parzystego stopnia całkowitego.
Generatory $x^2,y^2,xy$ spełniają jedyną relację
$UV=W^2$. Część stopnia $2k$ ma wymiar $2k+1$,
a część nieparzystego stopnia jest zerowa.
Dla $m=1$ relacja $W=UV$ eliminuje trzeci generator
i odzyskujemy cały pierścień $\mathbb C[x,y]$.
\end{przyklad}

\section{Grupy dwucykliczne}

\begin{twierdzenie}[Niezmienniki grupy dwucyklicznej]
\label{aa:tw:dicyclic-invariant-ring}
Niech $n\ge2$ i niech $\operatorname{Dic}_n$
działa na $\mathbb C^2$ przez macierze $A,J$
z Definicji~\ref{aa:def:binary-dicyclic}.
Połóżmy
\[
 u=(xy)^2,\quad v=x^{2n}+y^{2n},\quad
 w=xy(x^{2n}-y^{2n}).
\]
Wówczas
\[
 \mathbb C[x,y]^{\operatorname{Dic}_n}\simeq
 \mathbb C[U,V,W]/(W^2-UV^2+4U^{n+1}),
\]
gdzie $U,V,W$ przechodzą odpowiednio na $u,v,w$.
Szereg Hilberta wynosi
\[
 H_{\operatorname{Dic}_n}(z)
 =\frac{1+z^{2n+2}}{(1-z^4)(1-z^{2n})}.
\]
\end{twierdzenie}
\begin{proof}
Oznaczmy $t=xy$, $a=x^{2n}$, $b=y^{2n}$.
Z Twierdzenia~\ref{aa:tw:cyclic-invariant-ring}
niezmienniki podgrupy $\langle A\rangle$ mają bazę
$t^k$, $t^ka^j$, $t^kb^j$ ($k\ge0,j\ge1$),
a $ab=t^{2n}$. Działanie $J$ wysyła
$(t,a,b)$ na $(-t,b,a)$, ponieważ $J^{-1}(x,y)=(-y,x)$.
Ponieważ $A,J$ generują $\operatorname{Dic}_n$,
wystarczy wyznaczyć część pierścienia
$\mathbb C[x,y]^{\langle A\rangle}$ ustaloną przez $J$.
Jej bazę tworzą $t^{2k}$ oraz
\[
 t^{2k}(a^j+b^j),\qquad
 t^{2k+1}(a^j-b^j),\quad k\ge0,\ j\ge1.
\]
Istotnie dla każdego $k,j$ działanie na parze
$t^ka^j,t^kb^j$ zamienia wektory z czynnikiem $(-1)^k$;
przestrzeń ustalona jest jednowymiarowa.

Sumy $P_j=a^j+b^j$ spełniają rekurencję
$P_j=vP_{j-1}-u^nP_{j-2}$ dla $j\ge2$,
z $P_0=2,P_1=v$, bo $a+b=v$, $ab=u^n$.
Wielomiany $Q_j=\sum_{i=0}^{j-1}a^{j-1-i}b^i$
spełniają tę samą rekurencję dla $j\ge3$,
z $Q_1=1,Q_2=v$, i $a^j-b^j=(a-b)Q_j$.
Istotnie, dla $j\ge3$ iloczyn $(a+b)Q_{j-1}$
zawiera każdy skrajny monom $a^{j-1},b^{j-1}$
raz, a każdy monom pośredni dwa razy.
Odjęcie $abQ_{j-2}$ usuwa dokładnie jedno
wystąpienie każdego monomu pośredniego.
Każdy element wskazanej bazy jest więc kombinacją
wielomianów w $u,v,w$, a te trzy są niezmiennikami.
Ponadto
$w^2=t^2((a+b)^2-4ab)=uv^2-4u^{n+1}$.

Aby wykluczyć dalsze relacje, najpierw pokażemy
algebraiczną niezależność $u,v$.
Podstawienie $x=X,y=T/X$ w pierścieniu
$\mathbb C[T,X,X^{-1}]$ daje
$u=T^2$, $v=X^{2n}+T^{2n}X^{-2n}$.
Gdyby $\sum_{j=0}^r p_j(u)v^j=0$, $p_r\ne0$,
współczynnik przy najwyższej potędze $X^{2nr}$
wynosiłby $p_r(T^2)\ne0$, sprzeczność.
Każdy element ilorazu przez moniczną relację
kwadratową w $W$ zapisuje się jednoznacznie jako
$P(U,V)+WQ(U,V)$ (dzielenie przez wielomian moniczny).
Jeśli $P(u,v)+wQ(u,v)=0$, zamiana $x$ na $-x$
zachowuje $u,v$, a zmienia znak $w$.
Dodanie i odjęcie obu równości daje
$P(u,v)=0$ i $wQ(u,v)=0$.
Ponieważ $\mathbb C[x,y]$ nie ma dzielników zera,
$w\ne0$, a $u,v$ są algebraicznie niezależne,
otrzymujemy $P=Q=0$. To dowodzi iniektywności.

Monomy $u^iv^j$ i $wu^iv^j$ ($i,j\ge0$)
tworzą więc bazę jednorodną. Ich stopnie to
$4i+2nj$ oraz $2n+2+4i+2nj$, co daje szereg Hilberta.
\end{proof}

\begin{przyklad}[Niezmienniki grupy kwaternionowej]
Dla $n=2$ mamy $\operatorname{Dic}_2\simeq Q_8$.
Generatory $u,v$ mają stopień $4$, a $w$ stopień $6$;
jedyna relacja ma stopień $12$ i postać
$w^2=uv^2-4u^3$. W stopniu $4$ otrzymujemy bazę
$u,v$, w stopniu $6$ bazę $w$, a w stopniu $8$
bazę $u^2,uv,v^2$. Nie istnieją niezmienniki
dodatniego stopnia mniejszego niż $4$.
Natomiast $xy$ jest półniezmiennikiem stopnia $2$.
\end{przyklad}

\begin{wniosek}[Najmniejsze stopnie niezmienników]
\label{aa:wn:dicyclic-invariant-minimum}
Dla każdego $n\ge2$ minimalny dodatni stopień
niezmiennika grupy $\operatorname{Dic}_n$ wynosi $4$,
natomiast minimalny dodatni stopień półniezmiennika
wynosi $2$.
\end{wniosek}
\begin{proof}
Z szeregu Hilberta wynika brak niezmienników
w stopniach $1,2,3$ i istnienie $u$ stopnia $4$.
Wielomian $xy$ rozpina jednowymiarową podprzestrzeń
niezmienniczą, ale półniezmiennik stopnia $1$
dawałby niezmienniczą prostą w przestrzeni sprzężonej,
a więc jej jednowymiarowy anihilator byłby niezmienniczy
w $\mathbb C^2$. Jest to niemożliwe z nieprzywiedlności
w Wniosku~\ref{aa:wn:binary-primitive}.
\end{proof}

Prezentacje dla grup wyjątkowych $2T,2O,2I$ wymagają
dodatkowych argumentów o generatorach i relacjach.
Sam wzór Moliena nie zastępuje tych argumentów.
